\documentclass[12pt]{article}
\usepackage[T1]{fontenc}
\usepackage[utf8]{inputenc}
\usepackage{lmodern}
\usepackage{textcomp}
\usepackage{enumitem}
\usepackage{geometry}
\usepackage{graphicx}
\usepackage{amsmath, amssymb}
\geometry{margin=1in}
\begin{document}
\begin{center}
Economics - Undergraduate Level Assessment 13 \
Total Marks: 40 \
Answer all questions.
\end{center}

\section*{Multiple Choice (10 marks)}
\begin{enumerate}[label=\arabic*.]
\item Why is it difficult sometimes to definemonopolistically competitive industries?
\begin{enumerate}[label=(\alph*)]
    \item Product differentiation and loss of perfect substitutability
    \item Barriers to entry and exit in the market
    \item High consumer loyalty to a single brand
    \item Lack of government regulation
    \item Constant production costs regardless of output levels
\end{enumerate}
\item Common stock or bonds are considered more liquid assets thanan automobile. What would be the reason for this?
\begin{enumerate}[label=(\alph*)]
    \item Automobiles depreciate faster
    \item Stocks and bonds are physically smaller
    \item Automobiles are more expensive
    \item Highly efficient and fast markets exist for stocks and bonds
\end{enumerate}
\item Suppose a price floor is installed in the market for coffee. One result of this policy would be
\begin{enumerate}[label=(\alph*)]
    \item a decrease in the price of coffee.
    \item an increase in consumer surplus due to lower coffee prices.
    \item a persistent shortage of coffee in the market.
    \item a decrease in the demand for coffee.
    \item a decrease in consumer surplus due to higher coffee prices.
\end{enumerate}
\item A consequence of a price floor is
\begin{enumerate}[label=(\alph*)]
    \item an elimination of consumer surplus.
    \item a decrease in the demand for the good.
    \item a decrease in producer surplus.
    \item a persistent surplus of the good.
    \item an increase in total welfare.
\end{enumerate}
\item Households demand more money as an asset when
\begin{enumerate}[label=(\alph*)]
    \item the demand for goods and services decreases.
    \item the stock market crashes.
    \item bond prices fall.
    \item the nominal interest rate falls.
    \item the unemployment rate rises.
\end{enumerate}
\end{enumerate}

\section*{True / False (10 marks)}
\textit{Answer True or False}
\begin{enumerate}[label=\arabic*.]
\setcounter{enumi}{5}
\item True or False: An ARMA(p,q) model has an autocorrelation function that declines geometrically and a partial autocorrelation function that is zero after p lags.
\item True or False: The kinked demand curve consists of two distinct segments rather than being a single continuous function.
\item True or False: When price discrimination occurs, consumer surplus decreases as firms capture more of the consumer's willingness to pay.
\item True or False: Investment demand most likely increases when the stock market experiences a downturn, as investors seek safer assets.
\item True or False: Required reserves must be kept in a bank's vault and cannot be held in accounts at the Federal Reserve.
\end{enumerate}

\section*{Long Form Response (20 marks)}
\textit{Answer each question with a clear and complete explanation.}
\begin{enumerate}[label=\arabic*.]
\setcounter{enumi}{10}
\item Discuss the goals and economic rationale of the single-tax movement that emerged from the Ricardian theory of land rents in the late nineteenth century.
\item Discuss the relationship between the demand for stocks and bonds and the asset demand for money, analyzing how changes in one market can impact the other.
\end{enumerate}

\vfill
\noindent\textit{End of Paper}
\end{document}